\documentclass[11pt]{article}
\usepackage[margin=1in]{geometry}
\usepackage{booktabs}
\usepackage{amsmath}
\usepackage{amssymb}
\usepackage{xcolor}
\usepackage{hyperref}
\usepackage{array}
\hypersetup{colorlinks=true, linkcolor=blue!50!black, urlcolor=blue!50!black}

\newcommand{\good}[1]{\textcolor{green!45!black}{#1}}
\newcommand{\bad}[1]{\textcolor{red!70!black}{#1}}

\title{Can DINOv2 features score SD3.5 candidate images?\\
\large Four experiments on the ORCA visual target}
\author{}
\date{19 August 2026}

\begin{document}
\maketitle

\begin{abstract}
\noindent
We ran four experiments. The first rebuilds the ORCA paper's DINOv2 compression and matches its
reported number. The second shows that SD3.5's internal state can predict those compressed DINOv2
features. The third and fourth ask whether a DINOv2-based score can pick the better of eight images
generated from the same prompt. It cannot: a score that reads the prompt does no better than a
number that ignores the prompt. One category, \texttt{3d\_spatial}, is an exception.
\end{abstract}

\section{What we were testing and why}

ORCA trains a diffusion model with an extra loss. That loss pushes the model's internal state
towards a 64-number summary of DINOv2 features, through a matrix chosen by the prompt. We wanted to
know whether the same machinery can be used a different way: as a \emph{scorer} that ranks several
images generated from one prompt, without decoding them.

ORCA's predictor does not output a score. It outputs a 64-number vector,
\begin{equation}
\hat z(x_t, y) \;=\; K(\Delta_y)^\top h_t(x_t, y) \;\in\; \mathbb{R}^{64},
\end{equation}
where $h_t$ is the diffusion model's hidden state and $K(\Delta_y)$ is a matrix built from the
prompt. To get a score you must compare $\hat z$ against a target, or measure its size. So before
building anything, three things had to be true:

\begin{enumerate}
  \item the 64-number DINOv2 summary can be rebuilt (Experiment 1);
  \item SD3.5's hidden state can predict it (Experiment 2);
  \item DINOv2 features, given the prompt, tell you which image is better (Experiments 3 and 4).
\end{enumerate}

Item 3 is the one that matters. Items 1 and 2 only say the parts exist.

\section{Experiment 1: rebuild the DINOv2 compression}

\paragraph{Purpose.} Check we can reproduce ORCA's visual target before relying on it.

\paragraph{Method.} Take 5{,}000 images from COCO \texttt{train2017}. Centre-crop each to a square,
resize to $224\times224$ with bicubic interpolation, no flipping. Run DINOv2-L/14. Keep the patch
tokens after the final LayerNorm and drop the CLS token, giving $224/14 = 16$, so $16\times16 = 256$
vectors of length 1024 per image, and $5000 \times 256 = 1{,}280{,}000$ in total. Sample 100{,}000 of
them. Compute the mean $\bar v$ and run PCA in double precision. Keep the top 64 directions $P_{64}$.
The target for an image is
\begin{equation}
z_q(x) \;=\; P_{64}\bigl(v_q(x) - \bar v\bigr) \;\in\; \mathbb{R}^{64},
\qquad q = 1 \dots 256 .
\end{equation}

Two choices matter. We do \textbf{not} divide the tokens by their length, and we do \textbf{not}
whiten. ORCA's loss is a squared error over directions ordered by variance, so both steps would
change every eigenvalue.

\paragraph{Result.}

\begin{table}[h]
\centering
\caption{How much of the variation the top directions keep.}
\begin{tabular}{@{}rrrrrrrrr@{}}
\toprule
directions kept & 1 & 2 & 4 & 8 & 16 & 32 & \textbf{64} & 128 \\
\midrule
variation kept & 2.92\% & 5.25\% & 9.36\% & 14.65\% & 20.49\% & 27.60\% & \good{\textbf{37.05\%}} & 49.65\% \\
\bottomrule
\end{tabular}
\end{table}

\noindent
The ORCA paper reports \textbf{37.3\%} at 64 directions. We get \textbf{37.05\%}, a gap of
0.25 percentage points. The rebuild is correct.

\paragraph{Two things the paper claims that our numbers do not support.}

The paper assumes the eigenvalues shrink as $\lambda_i \le C\,i^{-\alpha}$ with $\alpha > 1$, and
uses that to bound how much is lost past rank $n$ by $C/((\alpha-1)n^{\alpha-1})$. We fitted
$\alpha$ over several ranges.

\begin{table}[h]
\centering
\caption{The fitted exponent $\alpha$ depends strongly on which components you fit over.}
\begin{tabular}{@{}lrl@{}}
\toprule
components fitted & $\alpha$ & \\
\midrule
1--64   & 0.676 & \bad{below 1} \\
1--128  & 0.649 & \bad{below 1} \\
1--256  & 0.648 & \bad{below 1} \\
1--512  & 0.698 & \bad{below 1} \\
64--1024  & 1.012 & just above 1 \\
128--1024 & 1.117 & above 1 \\
\bottomrule
\end{tabular}
\end{table}

\noindent
$\alpha$ is below 1 across the whole head of the spectrum. It only passes 1 far out in the tail, and
then only reaches 1.012, so the divisor $\alpha - 1 = 0.012$ makes the bound useless at rank 64.

Second, the paper expects performance to level off where the kept variation levels off. It does not
level off here: 37.05\% at 64, 49.65\% at 128, 65.66\% at 256, still climbing. 62.95\% of the
variation is still missing at rank 64. The paper's own Table~2 agrees with our reading, since its
GenEval score keeps rising to rank 128 and rank 64 was picked on FID instead.

\paragraph{Conclusion.} The compression rebuilds correctly. The paper's reason for choosing 64
directions does not hold up on our copy of its spectrum. This does not affect ORCA's reported FID
and GenEval gains.

\section{Experiment 2: can SD3.5's hidden state predict those 64 numbers?}

\paragraph{Purpose.} ORCA's predictor is a restricted linear map. Before building it, check that an
\emph{unrestricted} linear map can do the job. If the unrestricted one fails, the restricted one
must fail too.

\paragraph{Method.} Take 2{,}500 COCO images, none of them among the 5{,}000 used for the PCA. For
each image: encode to the SD3.5 VAE latent $x_0$, add noise at a fixed level,
\begin{equation}
x_\sigma \;=\; (1-\sigma)\,x_0 + \sigma\,\epsilon, \qquad \sigma = 0.5,
\end{equation}
run the frozen SD3.5 transformer, and capture the image tokens leaving block 8 of 24. Average-pool
those tokens to a $16\times16$ grid so they line up with the DINOv2 grid. Fit ridge regression from
the block-8 token to that patch's 64 DINOv2 numbers, on 2{,}000 images, and test on the held-out 500.
Repeat at three generation sizes.

\paragraph{Two controls, and why the first version was wrong.}

Our first version shuffled tokens against each other. That is the wrong floor. It destroys the grid
position as well as the image, so it lands near $-1$ whatever the model learned, and ``better than
shuffled'' is then true by construction. It reported $R^2 = 0.40$--$0.47$, which we withdrew.

SD3.5 adds position information to each token, and DINOv2 features have a position pattern (sky at
the top, object in the middle). So a ridge could score well by learning the average pattern per grid
slot, with no knowledge of the actual image. The two correct controls are:

\begin{itemize}
  \item \textbf{position only}: predict every image with the average DINOv2 map per grid slot,
        computed on the training images;
  \item \textbf{shuffled images}: predict image $i$'s map using image $j$'s features, keeping each
        token in its own grid slot.
\end{itemize}

\paragraph{Result.}

\begin{table}[h]
\centering
\caption{Held-out $R^2$ for predicting the 64 DINOv2 numbers from block 8. The column that matters
is the gain over position.}
\begin{tabular}{@{}lrrrr@{}}
\toprule
generation size & raw $R^2$ & position only & \textbf{gain over position} & shuffled images \\
\midrule
256 px  & 0.4029 & 0.1956 & \good{\textbf{0.2074}} & $-0.0399$ \\
512 px  & 0.4696 & 0.1956 & \good{\textbf{0.2741}} & $-0.1061$ \\
1024 px & 0.4729 & 0.1956 & \good{\textbf{0.2773}} & $-0.1049$ \\
\bottomrule
\end{tabular}
\end{table}

\begin{table}[h]
\centering
\caption{Other measurements from the same runs.}
\begin{tabular}{@{}lrrr@{}}
\toprule
generation size & cosine & feature length & effective rank \\
\midrule
256 px  & 0.6346 & 123.42 & 550.4 \\
512 px  & 0.6854 & 114.52 & 460.6 \\
1024 px & 0.6876 & 102.41 & 308.7 \\
\bottomrule
\end{tabular}
\end{table}

\paragraph{Conclusion.} Block 8 does carry image-specific DINOv2 information, worth $0.21$ to $0.28$
$R^2$ once the position pattern is removed. Position alone accounts for $0.196$, which is 42\% of the
raw 1024 px number, so the first version overstated the effect by nearly half.

512 px is the efficient choice: it gains $0.2741$ against 1024 px's $0.2773$, at a quarter of the
pixels. One caveat: at 256 px the SD3.5 grid is already $16\times16$, so pooling does nothing, while
512 px averages $2\times2$ tokens and 1024 px averages $4\times4$. The comparison across sizes is
therefore mixed up with how much averaging each one gets.

\section{The images and the labels used for scoring}

Experiments 3 and 4 need images and a ground truth. Both already existed.

\begin{table}[h]
\centering
\caption{The candidate pool. Every image was generated at SD3.5-Medium's own recommended settings:
1024 px, 40 steps, guidance 4.5, skip layers 7--9.}
\begin{tabular}{@{}lrrr@{}}
\toprule
prompt set & prompts & images per prompt & total images \\
\midrule
T2I-CompBench++ & 2{,}398 & 8 & 19{,}184 \\
GenEval2        &    800 & 8 &  6{,}400 \\
\midrule
total           & 3{,}198 & 8 & \textbf{25{,}584} \\
\bottomrule
\end{tabular}
\end{table}

We then scored \emph{every} candidate with the official evaluators, not just the selected one. This
took eight jobs, one per candidate index. The point is to have a ground truth that is not VQAScore,
so VQAScore can be a fair comparison rather than the answer key.

\begin{table}[h]
\centering
\caption{Average score per candidate index. They are nearly identical, as they must be: the eight
candidates are interchangeable random draws.}
\begin{tabular}{@{}lrrrrrrrr@{}}
\toprule
candidate & 0 & 1 & 2 & 3 & 4 & 5 & 6 & 7 \\
\midrule
CompBench & 0.4891 & 0.4896 & 0.4932 & 0.4871 & 0.5007 & 0.4895 & 0.4915 & 0.4888 \\
GenEval2  & 0.1791 & 0.1759 & 0.1765 & 0.1917 & 0.1816 & 0.1821 & 0.1907 & 0.1912 \\
\bottomrule
\end{tabular}
\end{table}

\subsection{How much is there to gain by picking well?}

\begin{table}[h]
\centering
\caption{Room to improve by choosing the best of the eight images, per prompt.}
\begin{tabular}{@{}lr@{}}
\toprule
best minus average, per prompt & \textbf{$+0.1429$} \\
best minus worst, per prompt   & $+0.3067$ \\
as a share of the average score & 45.1\% \\
variation within a prompt vs.\ between candidate indices & $1449\times$ \\
prompts where all eight images score the same & 7.0\% \\
\bottomrule
\end{tabular}
\end{table}

\begin{table}[h]
\centering
\caption{Room to improve, by category. The categories with the most room are the ones where you can
write a same-word wrong-binding version of the prompt.}
\begin{tabular}{@{}lrr@{}}
\toprule
category & room to improve & average score \\
\midrule
spatial      & 0.2191 & 0.2873 \\
numeracy     & 0.2021 & 0.5973 \\
3d\_spatial  & 0.1978 & 0.3785 \\
shape        & 0.1688 & 0.5570 \\
color        & 0.1376 & 0.7532 \\
texture      & 0.1333 & 0.6862 \\
complex      & 0.0689 & 0.3552 \\
non\_spatial & 0.0164 & 0.3147 \\
\bottomrule
\end{tabular}
\end{table}

\noindent
The room to improve, $+0.1429$, is larger than any training effect we have measured in this project.
Almost all of the variation is \emph{within} a prompt, which is what a scorer would exploit.

\subsection{167 prompts have no ranking to learn}

167 of the 2{,}398 prompts give all eight images the same score, so there is no order to predict.

\begin{table}[h]
\centering
\caption{Where the 167 flat prompts come from.}
\begin{tabular}{@{}lrrr@{}}
\toprule
category & prompts & flat & of which all-zero \\
\midrule
spatial     & 298 & 136 & \multirow{3}{*}{139 all-zero, 28 all non-zero} \\
numeracy    & 300 &  30 & \\
3d\_spatial & 300 &   1 & \\
others      & 300 each & 0 & \\
\bottomrule
\end{tabular}
\end{table}

\noindent
139 of them score zero on every image, meaning all eight failed. These prompts have no room to
improve, so they must stay in the denominator when averaging any selection gain. They are dropped
only from rank correlations, where they are undefined. One effect of this: any \texttt{spatial}
correlation below is measured on the 162 of 298 prompts where the evaluator separates the images at
all.

\section{Experiment 3: how good is VQAScore, the current best scorer?}

\paragraph{Purpose.} Set the bar. If a DINOv2 score is to be useful it has to be compared against
what is already available.

\paragraph{Method.} For each prompt, rank the eight images by VQAScore and by the official label, and
compute Kendall's $\tau_b$ between the two orders. Average over prompts. We do this \emph{within}
each prompt on purpose. Pooling all 19{,}184 images together would mostly measure which prompts are
easy, not which image is better.

\paragraph{Result.} Mean $\tau_b = \mathbf{+0.1536}$, 95\% interval $[+0.1389, +0.1670]$, over 2{,}231
prompts. \textbf{35.8\%} of prompts have $\tau \le 0$, meaning VQAScore orders them no better than
chance.

\begin{table}[h]
\centering
\caption{VQAScore against the official labels, by category.}
\begin{tabular}{@{}lrr@{}}
\toprule
category & mean $\tau_b$ & prompts \\
\midrule
spatial      & $+0.2674$ & 162 \\
color        & $+0.2034$ & 300 \\
shape        & $+0.1940$ & 300 \\
texture      & $+0.1878$ & 300 \\
numeracy     & $+0.1750$ & 270 \\
complex      & $+0.1386$ & 300 \\
non\_spatial & $+0.0788$ & 300 \\
3d\_spatial  & \bad{$+0.0377$} & 299 \\
\bottomrule
\end{tabular}
\end{table}

\paragraph{Conclusion.} VQAScore is a moderate scorer, not a strong one. It is close to useless on
\texttt{3d\_spatial} ($+0.0377$), which is also a category with a lot of room to improve ($0.1978$).
That gap is where a new scorer could earn its place.

\section{Experiment 4: can DINOv2 features rank the images?}

\paragraph{Purpose.} The main question. ORCA's target is a 64-number DINOv2 summary. Does that
summary, read with knowledge of the prompt, say which of eight images is better?

\paragraph{Method.} For each of the 25{,}584 images we stored the $16\times16\times64$ DINOv2 summary
and the CLS token. We then trained a small scorer,
\begin{equation}
S(x, y) \;=\; \operatorname*{reduce}_{q} \Bigl( \hat z_q(x)^\top W \, \hat e(y) \Bigr) + b ,
\end{equation}
where $\hat z_q$ is the length-normalised 64-number summary of patch $q$, $\hat e$ is the
length-normalised CLIP-L text embedding of the prompt, and $W$ is a single $64 \times 768$ matrix
(49{,}152 numbers). The reduce step is either the average over patches or the maximum. We train it to
order pairs of images correctly, using a hinge loss on pairs whose labels differ by at least half of
that category's typical within-prompt spread.

\paragraph{Splits.} 839 prompts to train, 359 to choose when to stop, 1{,}200 sealed for the final
number. The split is by prompt, so no test prompt is seen during training.

\paragraph{Controls.} Four of them.

\begin{enumerate}
  \item \textbf{Numbers that ignore the prompt.} If one of these matches the scorer, the prompt is
        doing no work.
  \item \textbf{Wrong prompt.} Feed a different prompt's text, averaged over 50 shuffles where no
        prompt keeps its own text.
  \item \textbf{Untrained matrix.} The same architecture with random $W$. Needed because this scorer
        reaches $\tau = +0.69$ on the training prompts even when fed pure noise, so it has ample
        capacity to memorise.
  \item \textbf{Per-category label spread} sets the pair margin, because the eight evaluators are on
        different scales.
\end{enumerate}

\begin{table}[h]
\centering
\caption{Typical within-prompt spread of the official label. A single fixed margin across these
would have starved \texttt{non\_spatial} of training pairs.}
\begin{tabular}{@{}lrrrrrrrr@{}}
\toprule
category & spatial & numeracy & shape & color & texture & 3d\_spatial & complex & non\_spatial \\
\midrule
spread & 0.2973 & 0.1652 & 0.1304 & 0.1201 & 0.1201 & 0.1175 & 0.0481 & 0.0108 \\
\bottomrule
\end{tabular}
\end{table}

\paragraph{Result.}

\begin{table}[h]
\centering
\caption{Mean within-prompt $\tau_b$ against the official labels, on 1{,}118 sealed prompts. The row
that decides the question is the last one.}
\begin{tabular}{@{}lrr@{}}
\toprule
 & average over patches & maximum over patches \\
\midrule
scorer that reads the prompt & $+0.0538$ & $+0.0430$ \\
\quad 95\% interval & $[+0.0356,+0.0719]$ & $[+0.0252,+0.0608]$ \\
\midrule
wrong prompt fed in            & $+0.0379$ & $+0.0259$ \\
untrained matrix               & $+0.0160$ & $+0.0203$ \\
\midrule
\textbf{best number that never reads the prompt} & \textbf{+0.0423} & \textbf{+0.0423} \\
\midrule
scorer minus prompt-free number & \bad{$+0.0115$} & \bad{$+0.0008$} \\
scorer minus wrong prompt       & $+0.0159$ & $+0.0172$ \\
\bottomrule
\end{tabular}
\end{table}

\begin{table}[h]
\centering
\caption{The numbers that never read the prompt. All are computed from the same 64-number DINOv2
summary.}
\begin{tabular}{@{}lrr@{}}
\toprule
number & mean $\tau_b$ & 95\% interval \\
\midrule
spread of the summary across image quadrants & $+0.0423$ & $[+0.0238,+0.0618]$ \\
how much the summary varies between patches  & $+0.0386$ & $[+0.0198,+0.0573]$ \\
average size of the first direction          & $+0.0337$ & $[+0.0155,+0.0521]$ \\
total size of all 64 directions              & $+0.0176$ & $[-0.0005,+0.0367]$ \\
length of the CLS token                      & $-0.0293$ & $[-0.0475,-0.0113]$ \\
\bottomrule
\end{tabular}
\end{table}

\begin{table}[h]
\centering
\caption{The scorer by category, both reduce steps. Bold marks intervals that stay above zero.}
\begin{tabular}{@{}lrrr@{}}
\toprule
category & average over patches & maximum over patches & prompts used \\
\midrule
3d\_spatial  & \good{$\mathbf{+0.1138}$ $[+0.0638,+0.1626]$} & \good{$\mathbf{+0.0719}$ $[+0.0267,+0.1173]$} & 150/150 \\
spatial      & \good{$\mathbf{+0.0748}$ $[+0.0045,+0.1452]$} & $+0.0356$ $[-0.0281,+0.0993]$ & 83/150 \\
numeracy     & \good{$\mathbf{+0.0685}$ $[+0.0097,+0.1262]$} & $+0.0069$ $[-0.0507,+0.0663]$ & 135/150 \\
shape        & \good{$\mathbf{+0.0651}$ $[+0.0187,+0.1104]$} & \good{$\mathbf{+0.0681}$ $[+0.0184,+0.1205]$} & 150/150 \\
color        & $+0.0380$ $[-0.0100,+0.0880]$ & $+0.0241$ $[-0.0221,+0.0714]$ & 150/150 \\
non\_spatial & $+0.0362$ $[-0.0142,+0.0865]$ & $+0.0217$ $[-0.0285,+0.0726]$ & 150/150 \\
complex      & $+0.0252$ $[-0.0252,+0.0743]$ & \good{$\mathbf{+0.0781}$ $[+0.0314,+0.1257]$} & 150/150 \\
texture      & $+0.0198$ $[-0.0293,+0.0681]$ & $+0.0310$ $[-0.0142,+0.0751]$ & 150/150 \\
\bottomrule
\end{tabular}
\end{table}

\paragraph{Conclusion.} The scorer fails. With the maximum reduce it beats a number that ignores the
prompt entirely by $+0.0008$, which is nothing. With the average reduce it wins by $+0.0115$. Feeding
it the wrong prompt costs only about $0.016$, so most of what little it does is not about the prompt.
For scale, VQAScore reaches $+0.1536$ on the same prompts and labels; this scorer reaches $+0.0538$.

One category breaks the pattern. On \texttt{3d\_spatial} the scorer reaches $+0.1138$
$[+0.0638, +0.1626]$ while VQAScore reaches $+0.0377$, about three times better, in a category with
$0.1978$ of room to improve. That result rests on 150 prompts.

\paragraph{Two things that weaken our own numbers.} We ran two reduce steps and are quoting the
better one, which flatters it; with intervals about $\pm 0.018$ wide, choosing between two options
after seeing the sealed set inflates the estimate. And the average reduce beating the maximum points
away from the story we expected: if the scorer worked by spotting the requested object in some patch,
the maximum should win.

\section{Summary}

\begin{table}[h]
\centering
\caption{All four experiments.}
\begin{tabular}{@{}p{3.6cm}p{4.4cm}p{5.6cm}@{}}
\toprule
experiment & result & verdict \\
\midrule
1. Rebuild the DINOv2 compression &
37.05\% kept at 64 directions, against the paper's 37.3\% &
\good{Works.} The paper's reason for choosing 64 does not hold up: $\alpha < 1$ across the head of
the spectrum, and the kept variation does not level off \\
\addlinespace
2. Predict the compression from SD3.5 block 8 &
$R^2$ gain over position: 0.2074 at 256 px, 0.2741 at 512 px, 0.2773 at 1024 px &
\good{Works.} 512 px is the efficient size. Position alone gives 0.1956, so our first number was
42\% inflated \\
\addlinespace
3. Measure VQAScore against official labels &
$\tau_b = +0.1536$; only $+0.0377$ on \texttt{3d\_spatial} &
Sets the bar. VQAScore is moderate, and weak on \texttt{3d\_spatial} \\
\addlinespace
4. Score candidates with DINOv2 features &
$\tau_b = +0.0538$, against $+0.0423$ for a number that ignores the prompt &
\bad{Fails.} Reading the prompt adds almost nothing. Exception: \texttt{3d\_spatial} at $+0.1138$ \\
\bottomrule
\end{tabular}
\end{table}

\subsection*{What we recommend}

\textbf{Do not train the full ORCA scoring kernel.} Its purpose is to rank candidates. The cheap
version of that test says DINOv2 features carry almost no prompt-specific ranking information. A
larger version of a signal that is not there will not find one.

\textbf{Look at \texttt{3d\_spatial} on its own.} It is the only place something appeared, and it is
also where VQAScore is worst, so it is the only place a new scorer would be worth having. Check it on
more than 150 prompts before trusting it.

\textbf{If anyone repeats this, keep the position control.} Our first version of Experiment 2 used a
token shuffle, which cannot fail, and reported $R^2 = 0.40$--$0.47$. Almost half of that was the
model knowing where each patch sits in the image.

\subsection*{Files}

\begin{table}[h]
\centering
\begin{tabular}{@{}ll@{}}
\toprule
what & where \\
\midrule
PCA fit & \texttt{phaseX/fit\_dino\_pca.py}, output \texttt{phaseX/dino\_pca/} \\
Block-8 prediction & \texttt{phaseX/block8\_readout.py}, output \texttt{phaseX/block8\_readout\_v2.json} \\
Candidate DINOv2 features & \texttt{phaseX/extract\_candidate\_dino.py}, output \texttt{phaseX/cand\_dino/} \\
Official per-candidate labels & \texttt{ablations/phaseX\_candeval.lsf}, output \texttt{phaseX/candeval/} \\
The scorer & \texttt{phaseX/dino\_probe.py}, output \texttt{phaseX/dino\_probe\_\{max,mean\}.json} \\
\bottomrule
\end{tabular}
\end{table}

\noindent
The PCA basis is pinned by two hashes: the image list \texttt{9ffea3f24ef55c96...} and the basis
itself \texttt{bf441fa9a739682d...}.

\end{document}
